\documentclass[12pt]{article}
\usepackage[utf8]{inputenc}
\usepackage[T1]{fontenc}
\usepackage{hyperref}
\title{Koro Zuba Language Resources for K-12 Curriculum Development}
\author{Compiled by Claude Sonnet 5 (claude-sonnet-5) for Andr\'as Kornai}
\date{}
\begin{document}
\maketitle

\section{Introduction}
Koro Zuba (vkz, Glottocode koro1324) is a Nupoid language of central Nigeria, spoken around Zuba in the Federal Capital Territory and in Niger State. One estimate gives 10,000--12,000 speakers (2020). It is also called Zuba Koro or Koro Ganagana. It has low lexical similarity with Koro Nulu, although speakers treat the two as one language.

\section{Primary Educational Resources}
None found.

\section{Linguistic and Reference Materials}
\begin{itemize}
\item Overview: \url{https://en.wikipedia.org/wiki/Koro_Zuba_language}
\item A 2021 sociolinguistic profile of Koro Nulu and Koro Zuba (Decker and colleagues), listed as a reference record: \url{https://glottolog.org/resource/reference/id/593564}
\end{itemize}

\section{Religious and Community Resources}
Joshua Project profile (reports positive language attitudes and a wish for language development): \url{https://joshuaproject.net/people_groups/20656}

\section{Digital and Audio Resources}
Global Recordings Network lists the language but reports no recordings: \url{https://globalrecordings.net/en/language/vkz}

\section{Curriculum Development Context}
The language is stable but not known to be taught in schools. Materials would have to be built from scratch with the community, starting from the sociolinguistic profile.
\end{document}